% res_scale.tex -- DRAFT 23 Oct 2026 (thesis result S): scalability. Updated 23 Oct afternoon with C, Q6 and Q1-Q5.
% Sources: card docs/experiments/2026-10-23_scalability.md (critic answered 22 Oct; bars FROZEN 23 Oct ~07:25 on Wesley's
% approval, decisions.md DECIDED 2026-10-23 (5): "not so important as long as it works"; eight quads reported "twin
% only, not real time", the real-time limit its own number); its "Results, C" (14:12, lab batch 17, R2700-R2711,
% results/analysis/s_c_twin_2026-10-23.json), "Results, Q6" (14:35, results/analysis/s_quality_bench_n8_2026-10-23.json)
% and "Results, Q1-Q5" (14:58, results/analysis/s_quality_replay_2026-10-23.json); decisions.md DECIDED 2026-10-23 S
% (Wesley ~14:33: N4's window opens 2 s after the last REJOIN_BLEND_OVER for the claim runs; the original window is
% reported beside it) and S's iteration count (~14:37: the 8-quad claims fly at 5 SQP iterations, s8_hover_R50_L100_it5
% and s8_chain4back_R50_L100_it5 x2 each, N1's comparator s4_hover_L100_it5 R2700/R2701; still "twin only, not real
% time"; the it1 8-quad runs stay exploratory). By the card's frozen rule (section 5 B and Q) the rig N_max is read at
% that count too: 5 iterations.
% Bench: registry S-BENCH-1022 (results/analysis/solve_scaling_tuf_2026-10-22.json), S-CALIB-1022
% (s_calib_replay_2026-10-22.json, s_calib_bench_n3_resin/_fast_2026-10-22.json, s_calib_replay_cores_2026-10-22.json;
% rig flights RIG-1022-h3 and -f3h). The bench rows are BEFORE the two speed-ups (plan_fast_refs), the residual inside the
% solve, and no 4 cm pivot in the reference extraction (the pivot adds about 4 ms at n = 3: whole tick p95 32.1 against
% 27.9 ms; 3.5 x 32.1 = 112 ms). Rig-scaled limit: 3.5 x whole tick p95 <= 100 ms, i.e. idle p95 <= 28.6 ms: n = 3 at
% 5 / 2 iterations (27.9 / 18.3 ms; n = 4 48.8 / 29.9), n = 4 at 1 (23.4; n = 5 40.7). Idle limit on latency p95 <= 60
% ms: 4 / 5 / 6.
% Eight-quad exploratory runs (1 iteration): R2464 (s8_hover), R2587 (s8_chain4), R2588 (s8_chain4back before the rejoin
% fixes), R2597 (s8_chain4back on the final rejoin tree).
% Not yet run: bar B (the bench with the speed-ups and the pivot, rig-scaled); the claim cells at 5 iterations (lab batch
% 18: s8_hover_R50_L100_it5 x2, s8_chain4back_R50_L100_it5 x2). NOT written up (Wesley has not decided): the card's
% PROPOSED "More complex cases" addendum (23 Oct ~15:05).
% Open, not in the text: two iterations on the rig (n <= 4 converged by every Q bar) would need a twin closed-loop check
% at it2, the reviewer and Wesley's word (a launch default).

\subsection{Scalability}
\label{sec:res_scale}

This section asks how many UAVs our planner can fly in real time on the rig's laptop, at how few solver iterations
per tick, and whether a larger team can carry a ring down a chain of unannounced releases and back. It gives three
answers apart. The real-time limit comes from a bench calibrated against the rig's own flights. How few iterations
are enough comes from the plans themselves, replayed against converged ones, and from the four-UAV twin flown at one
iteration and at five. The eight-UAV chain is a twin result, flown at a fraction of real time, and says nothing about
real time.

\subsubsection{The real-time limit on the rig laptop}

The planner's tick is one callback of its 100\,ms timer: the solve, the extraction of every UAV's references, and the
dynamics residual that the publish gate reads (Section~\ref{sec:controller}); a tick longer than the period delays
the next. A bench times that tick for a team of $n$ UAVs on the idle laptop, warm-started from its own previous plan
on a figure-eight at 0.2\,m/s and 1\,m cables, for 5, 2 and 1 sequential-quadratic-programming iterations per tick.

The bench is not the rig. Two rig flights of 22 October (three UAVs, five iterations), replayed tick by tick on the
same laptop with nothing else running, took the bench's time; the same ticks in flight took two to three and a half
times as long (Table~\ref{tab:scale_calib}). The factor is therefore the laptop in flight, shared with the rig's stack,
three trackers, motion capture and the recording, and not the flights' states. The limit uses the larger flight's
factor on the 95th percentile, rounded up to 3.5. It is a lower bound for a larger team, which would run more
trackers. One lead, untested on the rig: the laptop has performance and efficiency cores, the replay pinned to an
efficiency core took two and a half times as long as on a performance core, and in flight no tick was as fast as the
performance core's; the planner probably ran partly on the efficiency cores, and the next rig day pins it to the
performance cores.

\begin{table}[htbp]
\centering
\footnotesize
\caption[The bench calibrated against flight]{The planner's solve time on the rig laptop in flight against the same
laptop idle: three UAVs, five iterations, the residual inside the solve; warm ticks only. Replay: the flights' own
ticks solved again on the idle laptop. Bench: the timing bench at three UAVs. The rig factor is the flights' time over
the bench's; the limit of Table~\ref{tab:scale_nmax} uses 3.5 on the 95th percentile (2.1 on the median).}
\label{tab:scale_calib}
\begin{tabularx}{\textwidth}{>{\raggedright\arraybackslash}Xccc}
\toprule
 & Ticks & Median (ms) & 95th percentile (ms) \\
\midrule
Rig laptop in flight, hover & 264 & 36.7 & 67.8 \\
Rig laptop in flight, figure-eight & 573 & 40.4 & 69.4 \\
The same ticks replayed, idle (hover / figure-eight) & 264 / 573 & 18.5 / 19.2 & 19.4 / 20.6 \\
Bench, idle & 200 & 19.8 & 20.1 \\
\addlinespace
Rig factor, both flights pooled & 837 & 1.96 & 3.43 \\
\addlinespace
Hover replayed on one performance / one efficiency core & 264 & 18.4 / 47.7 & 18.9 / 48.9 \\
\midrule
Tick spacing in flight (both flights) & & -- & 116--117$^a$ \\
\bottomrule
\multicolumn{4}{p{0.95\textwidth}}{\footnotesize $^a$\,Measured near each tick's start: 10--12\,\% of ticks more than
110\,ms apart, none more than 150\,ms; the 10\,Hz loop lost no tick.}
\end{tabularx}
\end{table}
% S-CALIB-1022: h3 264 ticks 36.7/67.8 vs replay 18.5/19.4; f3h 573 ticks 40.4/69.4 vs 19.2/20.6; bench n = 3 19.8/20.1;
% k50/k95 1.86/3.37 (h3), 2.04/3.45 (f3h), pooled 1.96/3.43; pinned replay of h3 P-core 18.4/18.9, E-core 47.7/48.9.
% Lower bound: the rig's 60 ms budget cut held 11 (h3) and 37 (f3h) slow ticks to 3-4 iterations.

On that measure the rig laptop flies three UAVs in real time at five or two iterations a tick, and four at one
(Table~\ref{tab:scale_nmax}, Figure~\ref{fig:scale_bench}). The limit is read at five iterations, the count that the
checks of the next part pick: three UAVs. The calibration flights agree: with three UAVs at five iterations the rig's
ticks already came at about its period. The rig has flown four UAVs at five iterations (Section~\ref{sec:res_compute})
because its 60\,ms solve budget cuts the iterations of a slow tick; the limit here is for a tick that finishes its
iterations. The bench with the two speed-ups of the eight-UAV runs and the pivot, which gives the final number, is still
to run.

\begin{table}[htbp]
\centering
\footnotesize
\caption[The real-time limit]{The largest team the planner flies in real time, by solver iterations per tick. Idle:
the bench's latency (the solve and the references' extraction) under 60\,ms at the 95th percentile, the solver's own
throughput. Rig laptop: 3.5 times the bench's whole tick (with the residual) under the 100\,ms period at the 95th
percentile; the limit is read at five iterations. The bench ran before the two speed-ups of the eight-UAV runs and
without the 4\,cm pivot in the extraction, which adds about 4\,ms at three UAVs. No solve failed at any size.}
\label{tab:scale_nmax}
\begin{tabular}{lccc}
\toprule
Iterations per tick & 5 & 2 & 1 \\
\midrule
Idle laptop & 4 & 5 & 6 \\
Rig laptop & 3 & 3 & 4 \\
\addlinespace
Rig laptop, with the speed-ups and the pivot & [\dots] & [\dots] & [\dots] \\
\bottomrule
\end{tabular}
\end{table}
% Last row: the card's bar B (solve_scaling.py --n 2..10 --iters 5 2 1 --fast --pivot-dz 0.04 --rig-factor 3.5
% --tick-budget 100), not yet run (needs the laptop idle with no lab sync). That row is the thesis's number; the S-BENCH
% rows above are its baseline. By the card's rule (section 5: Q and C pick the count; Q6 missed i1 and i2) the rig N_max is
% read at 5 iterations: 3 on the old bench (the morning's "~4 quads" was the 1-iteration column).

\pendingfig{fig:scale_bench}{The planner's whole tick against team size on the idle rig laptop: the 95th percentile
of the solve, the references' extraction and the residual, warm-started, at 5, 2 and 1 iterations per tick (one line
each); dashed: 28.6\,ms, which the rig factor of 3.5 puts at the 100\,ms period in flight; the eight-UAV whole tick of
the twin's hover run, in wall time on the lab computer, for scale.}{the bench with the two speed-ups and the pivot}
% Data: results/analysis/solve_scaling_tuf_2026-10-22.json (now) and the bar-B bench JSON (to run); R2464's planner log
% for the twin point (tick_ms p95 101.2 ms wall). Tool: tools/solve_scaling.py output -> a new thesis_sim_figures
% function through the thesis-figure skill (log y axis; n = 2..9 or 10).

\subsubsection{How few iterations are enough}

Fewer iterations a tick leave the plan short of convergence. Two checks measure what that costs. Replayed tick by
tick, nine recorded flights, five from the twin and four from the rig with two to four UAVs, were solved again at one
and at two iterations from the iterate each tick started from, against the plan converged from the same iterate; an
untimed bench did the same at eight UAVs, with each warm start moved off the plan by about the rig's ring error
(Table~\ref{tab:scale_quality}). And the four-UAV twin flew a hover, the figure-eight at 0.2\,m/s and an unannounced
release at one iteration and at five, each from a floor start (Table~\ref{tab:scale_c}). The bars of both were written
before the runs.

At up to four UAVs, two iterations converge by every bar on all nine flights, the rig's among them. One iteration does
not: on every rig flight, and on the twin's line with a turn of the ring's yaw, its plans break the dynamics by more
than the bar allows. At eight UAVs one iteration misses every bar, and two miss only the thrust feed-forward's,
narrowly. Five pass every bar at eight, so by the rule written before the runs the eight-UAV claims fly at five
iterations, and the real-time limit above is read at five.

\begin{table}[htbp]
\centering
\footnotesize
\setlength{\tabcolsep}{2.5pt}
\caption[The plan at fewer iterations]{The plan at fewer iterations per tick against the plan converged from the same
iterate, replayed tick by tick: nine recorded flights (warm ticks only) and an untimed bench at eight UAVs on the
figure-eight at 0.2\,m/s, each warm start 1\,cm and 10\,cm/s off the plan. 95th percentiles of the dynamics residual (bar $10^{-3}$), of the error of the
references the trackers fly at the next node and the fifth (Ref., bars 5 and 10\,mm), and of the error of the thrust
feed-forward (FF, bar 0.3\,m/s$^2$). Bold: over a bar. The cost's ratio to the converged plan's (bars 1.10 on the median
and 1.50 on the 95th percentile, a tick over the residual bar counted as a miss) failed exactly where the residual did,
and no arm failed a solve. Five iterations at eight UAVs: $3.2 \times 10^{-7}$, 0.9 / 0.9\,mm and 0.12\,m/s$^2$.}
\label{tab:scale_quality}
\begin{tabular}{lcccccc}
\toprule
 & \multicolumn{3}{c}{Two iterations} & \multicolumn{3}{c}{One iteration} \\
\cmidrule(lr){2-4}\cmidrule(lr){5-7}
Flight (UAVs) & Residual & Ref.\ (mm) & FF (m/s$^2$) & Residual & Ref.\ (mm) & FF (m/s$^2$) \\
\midrule
Twin, hover (4) & $1.6{\times}10^{-7}$ & 0.2 / 0.2 & 0.023 & $1.9{\times}10^{-5}$ & 0.3 / 0.5 & 0.036 \\
Twin, figure-eight (4) & $1.3{\times}10^{-7}$ & 0.2 / 0.3 & 0.030 & $1.3{\times}10^{-5}$ & 0.4 / 0.7 & 0.045 \\
Twin, line, yaw turn (4) & $6.1{\times}10^{-6}$ & 0.5 / 0.7 & 0.062 & $\mathbf{2.3{\times}10^{-3}}$ & 1.2 / 2.6 & 0.107 \\
Twin, figure-eight (3) & $2.9{\times}10^{-7}$ & 0.3 / 0.3 & 0.033 & $2.0{\times}10^{-5}$ & 0.4 / 0.7 & 0.050 \\
Twin, figure-eight (2) & $3.1{\times}10^{-7}$ & 0.3 / 0.4 & 0.041 & $4.3{\times}10^{-5}$ & 0.6 / 0.9 & 0.063 \\
\addlinespace
Rig, hover (3) & $6.3{\times}10^{-5}$ & 0.8 / 2.2 & 0.098 & $\mathbf{3.0{\times}10^{-3}}$ & 1.6 / 4.1 & 0.154 \\
Rig, figure-eight (3) & $5.5{\times}10^{-5}$ & 1.3 / 3.0 & 0.144 & $\mathbf{2.3{\times}10^{-3}}$ & 2.3 / 5.6 & 0.227 \\
Rig, line, yaw turn (3) & $1.3{\times}10^{-4}$ & 1.5 / 3.3 & 0.165 & $\mathbf{6.6{\times}10^{-3}}$ & 3.1 / 6.8 & 0.279 \\
Rig, figure-eight (2) & $8.6{\times}10^{-5}$ & 2.2 / 4.6 & 0.247 & $\mathbf{7.4{\times}10^{-3}}$ & 3.7 / 10.0 & \textbf{0.378} \\
\addlinespace
Bench, figure-eight (8) & $8.8{\times}10^{-5}$ & 3.2 / 5.1 & \textbf{0.316} & $\mathbf{9.3{\times}10^{-3}}$ & \textbf{7.3} / \textbf{14.5} & \textbf{0.60} \\
\bottomrule
\end{tabular}
\end{table}
% Q1-Q5: twin R2348 (b4 hover), R2352 (b4 figure-8 size 1.0), R2356 (b4 line_spin, constant-rate turn), R2343 (b3
% figure-8), R2360 (b2 figure-8); rig RIG-1022-h3, -f3h, -ls3c, -f2slope. i1 ticks over the residual bar: R2356 117,
% h3 49, f3h 148, ls3c 232, f2slope 294; i2: f2slope 4. Q6 (n = 8, noise 0.01, untimed): i5 1.000/1.001, 3.2e-7,
% 0.9/0.9, 0.12; i2 1.003/1.010, 8.8e-5, 3.2/5.1, 0.316; i1 170 of 200 ticks over the residual bar. The bench flies
% solve_scaling's figure-8 (size 1.0, 0.2 m/s; --traj-radius 1.0) with a kinematic plant, --noise 0.01 (1 cm, 10 cm/s).

In the twin's closed loop one iteration a tick flew four UAVs as well as five in every cell: inside the carry bands of
the hover and the figure-eight, within the bars set from the five-iteration runs through the lift, the path and the
release, and with no published plan's residual over the bars' limit. It flew the figure-eight a little closer in both repeats, likely because its plan
reached the trackers younger, which was not tested. Its published plans did carry a residual, largest at the release,
that the planner's publish gate never reads at one iteration, since the gate reads it only when the solve budget cuts
a tick; it stayed under the closed loop's limit but over the replay's bar.

\begin{table}[htbp]
\centering
\footnotesize
\setlength{\tabcolsep}{2.4pt}
\caption[Four UAVs at five and at one iteration]{The four-UAV twin at five and at one iteration per tick: the 0.86\,kg
ring at 1.0\,m on 1\,m cables, floor start; hover and the figure-eight (size 1.0 at 0.2\,m/s after a 5\,s hold) on
plates 0, 3, 6 and 9; the release on the join's layout (plates 3, 7 and 11, the UAV on plate 9 released unannounced,
leaving a Y). Two runs each (a / b), one code tree. Path error: the clock-free path error (the hover from 2\,s after
the top to LAND); path tilt peak over the same window, which in the release cell holds the release; lift: from the ring's
reference 1\,cm above its start to 2\,s after the top; release: from the 5\,s mean before it. Residual: the largest
of a published plan, against the closed loop's limit of $10^{-2}$. In brackets, the bars for one iteration, set from
the two runs at five. No run aborted, dropped or failed a solve; the planner flew three UAVs 0.28--0.32\,s after the
release and the ring settled within 0.11$^\circ$ of level.}
\label{tab:scale_c}
\begin{tabular}{llccccc}
\toprule
 & & Path & \multicolumn{3}{c}{Tilt peak ($^\circ$)} & \\
\cmidrule(lr){4-6}
Cell & Iter. & error (mm) & Path & Lift & Release & Residual \\
\midrule
Hover & 5 & 1.34 / 1.32 & 0.25 / 0.20 & 1.77 / 1.87 & -- & $1.1{\times}10^{-8}$ \\
 & 1 & 1.35 / 1.37 [4.33] & 0.24 / 0.21 [1.23] & 1.82 / 1.87 [2.82] & -- & $1.5{\times}10^{-4}$ \\
\addlinespace
Figure-eight & 5 & 11.21 / 11.43 & 1.05 / 1.05 & 1.80 / 1.78 & -- & $1.2{\times}10^{-8}$ \\
 & 1 & 9.66 / 9.84 [14.32] & 0.89 / 0.93 [2.05] & 1.82 / 1.75 [2.79] & -- & $4.2{\times}10^{-4}$ \\
\addlinespace
Release & 5 & -- & 6.75 / 5.79 & 2.57 / 2.54 & 8.25 / 7.40 & $2.5{\times}10^{-7}$ \\
 & 1 & -- & 6.46 / 6.96 [7.27] & 2.63 / 2.66 [3.55] & 8.00 / 8.46 [8.82] & $2.3{\times}10^{-3}$ \\
\bottomrule
\end{tabular}
\end{table}
% R2700/R2701 | R2702/R2703 | R2704/R2705 | R2706/R2707 | R2708/R2709 | R2710/R2711 (lab batch 17, tree be34f95 + patch
% 537551ba, RTF 0.22-0.27, all 12 valid; C1-C7 hold in both repeats of every cell: C CONFIRMS). Closest margins R2711:
% C3 by 0.31 deg, C7's release peak by 0.37 deg. b22 bands: hover <= 22.0 mm / 3.21 deg, figure-8 <= 33.0 / 4.06.
% Published res_eq over 1e-3: 5 ticks in each it1 release run. The it5 pair is also plan_fast_refs' twin smoke run.

One iteration is therefore enough for four UAVs in the twin's closed loop, but its plans are not converged on the
rig's flights and fall short of every bar at eight; two converge for every fleet of up to four, the rig's flights
included.

\subsubsection{Eight UAVs in the twin}

The twin's large ring is 0.5\,m in radius and weighs 0.9\,kg, larger than the rig's ring and not much heavier, with
eight UAVs at 45$^\circ$ on 1\,m cables. The runs flown so far, all exploratory, flew one iteration a tick; the claim
runs fly five. Both fly two changes that shorten the tick without changing its values: every UAV's references come
from one mapped call, and the residual is computed after the references go out. The world runs at a capped fraction
of real time, so a tick that takes the planner its whole period in wall time takes a fraction of it in simulated time:
the result is about the plan and the physics, not about real time, and each run reports its achieved real-time factor
and its late ticks first. Table~\ref{tab:scale_twin} gives the four runs flown so far, one each.

Eight UAVs lifted the ring from the floor and hovered it level. In the chain, every other UAV was released
unannounced, in opposite pairs, so that each release left gaps of at most 90$^\circ$, down to a square of four; the
ring's tilt peaked at a few degrees and the drop rule never fired. In the return, the four freed UAVs, held at
staging points beyond their own plates, rejoined one at a time onto those plates (Section~\ref{sec:res_rejoin}), the
team grew back to eight and landed. Flown again on the final code tree of the rejoin, the chain gave the same
picture (Figure~\ref{fig:scale_chain}). Even in the twin, with the world slowed, the eight-UAV tick at one iteration
took about its whole period of wall time on a 12-core slot of the lab computer, half of it the residual: it is the
whole tick, not the solve, that would overrun first.

The claim runs, two hovers and two chains with the return at five iterations, are still to fly. Their bars were
frozen first, and one has been moved since, still before any claim run flew: the team back at eight is scored from
2\,s after the last rejoin's blend ends, because the rejoin's 10\,s blend, added after the bars were written, put the
last fold-in inside the window first written, from 2\,s after the resize to eight; that window is reported beside it.

\begin{table}[htbp]
\centering
\footnotesize
\setlength{\tabcolsep}{3pt}
\caption[Eight UAVs in the twin]{Eight UAVs in the rig twin, twin only and not real time: the 0.9\,kg ring of radius
0.5\,m at 1.0\,m, eight UAVs at 45$^\circ$ on 1\,m cables, floor start, hover, one iteration per tick; one run each,
exploratory. Releases unannounced, in opposite pairs, down to the square; each rejoin starts 20\,s after the previous
weld and holds the ring through a 10\,s fold-in blend. The first three runs flew before the rejoin's fixes
(Section~\ref{sec:res_rejoin}), the last on its final tree. Ring tilt over the airborne flight. $^a$\,From each
rejoin's start to its rod's end on the plate (before the fixes) or to the weld (final tree); the weld follows the
plate by about 0.1\,s. $^b$\,Ring position error (rms) and tilt peak with the team back at eight, from 2\,s after the
last blend ends to LAND, and in brackets from 2\,s after the resize to eight.}
\label{tab:scale_twin}
\begin{tabularx}{\textwidth}{>{\raggedright\arraybackslash}p{3.0cm}>{\raggedright\arraybackslash}X>{\raggedright\arraybackslash}X>{\centering\arraybackslash}p{1.9cm}c}
\toprule
Run & Releases & Rejoins$^a$ & Ring tilt, median / peak ($^\circ$) & Real-time factor \\
\midrule
Hover & -- & -- & 0.02 / 0.08 & 0.28 \\
8 $\to$ 4 & four, no drop or abort & -- & 0.13 / 3.48 & -- \\
8 $\to$ 4 $\to$ 8, before the fixes & four, no drop or abort & four, 16.1--18.4\,s & 0.18 / 3.25 & 0.2 \\
8 $\to$ 4 $\to$ 8, final tree & four, no drop or abort & four, 15.1--19.4\,s; each rod's end within 1.3--1.6\,cm of
  its plate through the blend; back at eight$^b$ 6.3\,mm, 0.89$^\circ$ (19.3\,mm, 1.65$^\circ$) & 0.52 / 3.35 & -- \\
\addlinespace
Claim cells, five iterations: hover; 8 $\to$ 4 $\to$ 8 & [two runs each, to fly] & [\dots] & [\dots] & [\dots] \\
\midrule
\multicolumn{5}{p{0.93\textwidth}}{Planner tick in the hover run, wall time on the lab computer: solve 44 / 60\,ms
(median / 95th percentile), whole tick 79 / 101\,ms, of which the residual 44--51\,ms; 0.5\,\% of ticks late in
simulated time.} \\
\bottomrule
\end{tabularx}
\end{table}
% R2464 | R2587 | R2588 (wall 28 min, RTF ~0.2) | R2597. Every run landed. R2597's back-at-eight (s_metrics n4 / n4_resize,
% decisions 23 Oct S (1)): revised 6.3 mm / 0.89 deg, original 19.3 mm / 1.65 deg; the bars the 1-iteration claim would
% have had were 5.72 mm / 1.44 deg (N1 from R2702/R2703); the claim runs at 5 iterations take N1 from R2700/R2701. The
% claim cells (N1-N6 of the card): s8_hover_R50_L100_it5 x2, s8_chain4back_R50_L100_it5 x2 (lab batch 18, --slots 2
% --cpus 12 --wall-cap-min 150), scored by tools/s_metrics.py (clock fix of 23 Oct 14:33).

\pendingfig{fig:scale_chain}{Eight UAVs to four and back in the rig twin (twin only, not real time): the 0.9\,kg ring
of radius 0.5\,m, hover, five iterations per tick. (a) The team size the planner flies, with each unannounced release
and each rejoin's start, weld and end of blend marked; (b) the ring's tilt, against the chain's bars of 15$^\circ$ at a
release and 10$^\circ$ at a rejoin; (c) plan views of the formation at eight, at four and back at eight.}{a claim run
of the eight-UAV chain with the return}
% Data: R2597 now (events.csv, run.csv, planner log n), the claim runs later. Tool: tools/s_metrics.py windows -> a new
% thesis_sim_figures function through the thesis-figure skill; panel (c) could be three frames of the run's video
% (results/videos, tools/render_video.py) instead.

\paragraph{Answer.} On the rig laptop our planner flies three UAVs in real time at five iterations a tick, the count
its quality checks pick, on a bench calibrated against the rig's own flights; the final number waits for the bench
with the speed-ups. Fewer iterations are not free: one iteration is enough for four UAVs in the twin's closed
loop, but its plans are not converged on the rig's flights and miss every bar at eight UAVs, while two converge
for every fleet of up to four. In the twin, with the clock taken away, eight UAVs at one iteration carried a larger
ring through four unannounced releases to a square and back to eight, with its tilt peaking at a few degrees, in single
exploratory runs; the claim runs, at five iterations, are to fly. The eight-UAV result is the twin's, and not real
time.
